\begin{abstract}
Developing a new Alzheimer's disease (AD) drug takes more than a decade, and almost all candidates fail. \emph{Drug repurposing}---finding already-approved drugs that might also help AD---is a faster route. A well-established way to screen for such drugs treats the cell's proteins as a graph: a drug is a promising candidate if the proteins it acts on sit unusually close, in shortest-path distance, to the proteins linked to AD. Deciding what ``unusually close'' means requires comparing each drug against thousands of random look-alikes, which makes the screen expensive. We propose a CUDA implementation of this screen. It computes all shortest-path distances on the human protein-interaction network with a bit-parallel breadth-first search, then scores every drug against thousands of random samples in parallel on the GPU. We compare its execution time against a sequential C++ version, check that both produce identical distances, and test whether the resulting ranking places drugs that have already entered AD clinical trials near the top. As an optional extension, we express the choice of a small drug combination as a Quadratic Unconstrained Binary Optimization (QUBO) problem and benchmark an exact GPU brute-force solver and classical heuristics against the Quantum Approximate Optimization Algorithm (QAOA), simulated on the GPU with NVIDIA cuQuantum (no quantum hardware is used).
\end{abstract}